% GENERATED FILE. Edit canonical lessons, then run npm run generate:books.
% canonical-source-hash: fnv1a64:229280255709b754
% canonical-lessons: ZH-C320-guaduan, ZH-C320-boda, ZH-C320-datong, ZH-C320-shuoqi, ZH-C320-yuyin

\chapter{On the Phone}
\label{ch:zh-on-the-phone}
\hlchaptermodality{320}

\begin{chapteropening}
\textbf{By the end of this chapter:} \emph{I can say someone hung up, dial a number, get through, bring something up and talk about a voice message.}

Complete the last lesson of chapter 320: I can say someone hung up, dial a number, get through, bring something up and talk about a voice message.
\end{chapteropening}

\section[guàduàn]{guàduàn — to hang up the phone; to end a call}

\label{lesson:ZH-C320-guaduan}



\begin{quote}
Before the new one: say the Mandarin for to write to each other, then the Mandarin for phone charges.
\end{quote}

\subsection*{You'll want to know: guàduàn}
\emph{guàduàn} — \textquotedblleft{}to hang up the phone; to end a call\textquotedblright{}.

\begin{quote}
\emph{tā hái méi shuō wán jiù guàduàn le diànhuà} — He hung up before he had finished.
\end{quote}

Its characters come later; for now it is a word for the ear.

\begin{cousinweb}
\emph{guàduàn} is built from \emph{guà} 'to hang' and \emph{duàn} 'to cut off', both inherited from Middle Chinese, \emph{kweaH} and \emph{dwanX} (Guangyun, in Baxter's transcription), which glosses the first 'to suspend' and the second 'to sever'.
\end{cousinweb}

\begin{grammarlens}[title={What you've built}]
One more word for telling what happened.
\end{grammarlens}

\subsection*{Your turn}
\begin{itemize}
\raggedright
  \item \emph{Say it:} \emph{guàduàn}
  \item \emph{Say it:} \emph{guàduàn}, once more
  \item \emph{Say it:} \emph{tā hái méi shuō wán jiù guàduàn le diànhuà} — He hung up before he had finished.
  \item \emph{Recall:} say the Mandarin for to write to each other, then the Mandarin for phone charges, then say \emph{guàduàn} again
\end{itemize}

\begin{tcolorbox}[breakable,colback=teal!4,colframe=teal!35!black,title={Before you move on}]
What does guàduàn mean? (\textquotedblleft{}To hang up the phone; to end a call\textquotedblright{}.) Say it once more.
\end{tcolorbox}
\section[bōdǎ]{bōdǎ — to dial a number}

\label{lesson:ZH-C320-boda}



\begin{quote}
Before the new one: say the Mandarin for phone charges, then the Mandarin for to end a call.
\end{quote}

\subsection*{You'll want to know: bōdǎ}
\emph{bōdǎ} — \textquotedblleft{}to dial a number\textquotedblright{}.

\begin{quote}
\emph{yǒu shì qǐng bōdǎ zhège hàomǎ} — If anything comes up, please call this number.
\end{quote}

Its characters come later; for now it is a word for the ear.

\begin{cousinweb}
\emph{bōdǎ} is built from \emph{bō} and \zh{打} 'to strike', both inherited from Middle Chinese, \emph{pat} and \emph{taengX} (Guangyun, in Baxter's transcription), which glosses the first 'to set in order' and 'to remove' and the second 'to hit'. The sense 'to dial' is modern. Middle Chinese \emph{taengX} would not regularly give dǎ.
\end{cousinweb}

\begin{grammarlens}[title={What you've built}]
One more word for telling what happened.
\end{grammarlens}

\subsection*{Your turn}
\begin{itemize}
\raggedright
  \item \emph{Say it:} \emph{bōdǎ}
  \item \emph{Say it:} \emph{bōdǎ}, once more
  \item \emph{Say it:} \emph{yǒu shì qǐng bōdǎ zhège hàomǎ} — If anything comes up, please call this number.
  \item \emph{Recall:} say the Mandarin for phone charges, then the Mandarin for to end a call, then say \emph{bōdǎ} again
\end{itemize}

\begin{tcolorbox}[breakable,colback=teal!4,colframe=teal!35!black,title={Before you move on}]
What does bōdǎ mean? (\textquotedblleft{}To dial a number\textquotedblright{}.) Say it once more.
\end{tcolorbox}
\section[dǎtōng]{dǎtōng — to get through on the phone}

\label{lesson:ZH-C320-datong}



\begin{quote}
Before the new one: say the Mandarin for to end a call, then the Mandarin for to dial a number.
\end{quote}

\subsection*{You'll want to know: dǎtōng}
\emph{dǎtōng} — \textquotedblleft{}to get through on the phone\textquotedblright{}.

\begin{quote}
\emph{dì'èr tiān wǒ cái dǎtōng tā de diànhuà} — I didn't get through to him until the next day.
\end{quote}

Its characters come later; for now it is a word for the ear.

\begin{cousinweb}
\emph{dǎtōng} is built from \zh{打} 'to strike' and \emph{tōng} 'to reach', both inherited from Middle Chinese, \emph{taengX} and \emph{thuwng} (Guangyun, in Baxter's transcription), which glosses the first 'to hit' and the second 'to reach'. Middle Chinese \emph{taengX} would not regularly give dǎ: the modern vowel is irregular.
\end{cousinweb}

\begin{grammarlens}[title={What you've built}]
One more word for telling what happened.
\end{grammarlens}

\subsection*{Your turn}
\begin{itemize}
\raggedright
  \item \emph{Say it:} \emph{dǎtōng}
  \item \emph{Say it:} \emph{dǎtōng}, once more
  \item \emph{Say it:} \emph{dì'èr tiān wǒ cái dǎtōng tā de diànhuà} — I didn't get through to him until the next day.
  \item \emph{Recall:} say the Mandarin for to end a call, then the Mandarin for to dial a number, then say \emph{dǎtōng} again
\end{itemize}

\begin{tcolorbox}[breakable,colback=teal!4,colframe=teal!35!black,title={Before you move on}]
What does dǎtōng mean? (\textquotedblleft{}To get through on the phone\textquotedblright{}.) Say it once more.
\end{tcolorbox}
\section[\texorpdfstring{\zh{说起} (shuōqǐ)}{说起 (shuōqǐ)}]{\zh{说起} (shuōqǐ) — to mention; to bring up a subject}

\label{lesson:ZH-C320-shuoqi}



\begin{quote}
Before the new one: say the Mandarin for to dial a number, then the Mandarin for to get through on the phone.
\end{quote}

\subsection*{You'll want to know: \zh{说起}}
\textbf{\zh{说起}} — \emph{shuōqǐ} — \textquotedblleft{}to mention; to bring up a subject\textquotedblright{}.

\begin{quote}
\textbf{\zh{他在信里说起了你}} \emph{tā zài xìn lǐ shuōqǐ le nǐ} — He mentioned you in his letter.
\end{quote}

\begin{cousinweb}
\textbf{\zh{说起}} (\emph{shuōqǐ}) is built from \zh{说} 'to say' and \zh{起} 'to rise', both inherited from Middle Chinese, \emph{sywet} and \emph{khiX} (Guangyun, in Baxter's transcription), which glosses the first 'to tell' and the second 'to rise' and 'to stand up'.
\end{cousinweb}

\begin{grammarlens}[title={What you've built}]
One more word for telling what happened.
\end{grammarlens}

\subsection*{Your turn}
\begin{itemize}
\raggedright
  \item \emph{Say it:} \emph{shuōqǐ}
  \item \emph{Say it:} \emph{shuōqǐ}, once more
  \item \emph{Say it:} \emph{tā zài xìn lǐ shuōqǐ le nǐ} — He mentioned you in his letter.
  \item \emph{Recall:} say the Mandarin for to dial a number, then the Mandarin for to get through on the phone, then say \emph{shuōqǐ} again
\end{itemize}

\begin{tcolorbox}[breakable,colback=teal!4,colframe=teal!35!black,title={Before you move on}]
What does \zh{说起} mean? (\textquotedblleft{}To mention; to bring up a subject\textquotedblright{}.) Say it once more.
\end{tcolorbox}
\section[\texorpdfstring{\zh{语音} (yǔyīn)}{语音 (yǔyīn)}]{\zh{语音} (yǔyīn) — voice message; speech sounds}

\label{lesson:ZH-C320-yuyin}



\begin{quote}
Before the new one: say the Mandarin for to get through on the phone, then the Mandarin for to bring up a subject.
\end{quote}

\subsection*{You'll want to know: \zh{语音}}
\textbf{\zh{语音}} — \emph{yǔyīn} — \textquotedblleft{}voice message; speech sounds\textquotedblright{}.

\begin{quote}
\emph{wǒ tīng le māma de yǔyīn} — I listened to my mum's voice message.
\end{quote}

\begin{cousinweb}
\textbf{\zh{语音}} (\emph{yǔyīn}) is built from \zh{语} 'speech' and \zh{音} 'a sound', both inherited from Middle Chinese, \emph{ngjoX} and \emph{'im} (Guangyun, in Baxter's transcription), which glosses the first 'to discuss' and the second 'a sound'. A recorded message is a modern sense.
\end{cousinweb}

\begin{grammarlens}[title={What you've built}]
One more word for telling what happened.
\end{grammarlens}

\subsection*{Your turn}
\begin{itemize}
\raggedright
  \item \emph{Say it:} \emph{yǔyīn}
  \item \emph{Say it:} \emph{yǔyīn}, once more
  \item \emph{Say it:} \emph{wǒ tīng le māma de yǔyīn} — I listened to my mum's voice message.
  \item \emph{Recall:} say the Mandarin for to get through on the phone, then the Mandarin for to bring up a subject, then say \emph{yǔyīn} again
\end{itemize}

\begin{tcolorbox}[breakable,colback=teal!4,colframe=teal!35!black,title={Before you move on}]
What does \zh{语音} mean? (\textquotedblleft{}Voice message; speech sounds\textquotedblright{}.) Say it once more.
\end{tcolorbox}
